\subsubsection{Compactness, Connectedness, and Completeness}

% Source pages provided: Rudin Ch.2, Def. 2.18 and examples around 2.21 for bounded/perfect terminology only.
% Sources mapped: Abbott Ch.3, Ch.8.2; Pugh Ch.2; Tao II Ch.1-2; Grinberg Ch.9-13.
% TODO(source pages unavailable): Add compactness, connectedness, completeness, Heine-Borel, and Baire-category material once source pages are provided.

This subsection is reserved for global structure: compactness, connectedness, completeness, and the theorems that make them useful. The provided pages only support a few metric-set prerequisites, so the main compactness and connectedness sequence remains source-gated.

\begin{definition} \label{def:bounded-and-perfect-metric-sets} \hlt{Bounded and Perfect Metric Sets}\\
Let $E$ be a subset of a metric space $X$.
\begin{enumerate}[label=(\roman*)]
\setlength{\itemsep}{0pt}
\item $E$ is \hlt{bounded} if there exist $M\in\R$ and $q\in X$ such that $d(p,q)<M$ for all $p\in E$.
\item $E$ is \hlt{perfect} if $E$ is closed and every point of $E$ is a limit point of $E$.
\end{enumerate}
The definitions of closed set and limit point are in Definition \ref{def:metric-topology-core-definitions}.
\end{definition}

\begin{definition} \label{def:compact-connected-complete-core-source-todo} \hlt{Core Structural Definitions to Source}\\
TODO(source pages unavailable): Add the sourced definitions of compactness, connectedness, Cauchy sequences, and completeness. Keep perfect-set terminology aligned with Definition \ref{def:bounded-and-perfect-metric-sets}.
\end{definition}

\begin{theorem} \label{thm:compact-connected-complete-main-source-todo} \hlt{Main Structural Theorems to Source}\\
TODO(source pages unavailable): Add the sourced theorem sequence linking compactness, connectedness, and completeness in metric or Euclidean settings.
\end{theorem}

\begin{method} \label{meth:compact-connected-complete-proof-patterns-source-todo} \hlt{Proof Patterns to Source}\\
TODO(source pages unavailable): Add proof commentary for cover arguments, sequence arguments, and separation arguments only after source coverage is available.
\end{method}
